% =====================================================================
%  Creating a Niagara N4 Module — Workbench Wizard + IntelliJ IDEA setup
%  Manual documenting the operator workflow (ColdRoomPan example).
%  Build: pdflatex main.tex  (run twice for the ToC / references)
% =====================================================================
\documentclass[11pt,a4paper]{article}

\usepackage[utf8]{inputenc}
\usepackage[T1]{fontenc}
\usepackage[margin=2.4cm]{geometry}
\usepackage{graphicx}
\usepackage{float}
\usepackage{booktabs}
\usepackage{enumitem}
\usepackage{xcolor}
\usepackage{listings}
\usepackage{caption}
\usepackage{url}
\usepackage[hidelinks]{hyperref}

\graphicspath{{img/}}

% --- palette / callout box -------------------------------------------
\definecolor{niagaraRed}{RGB}{197,42,42}
\definecolor{codeBg}{RGB}{245,245,245}
\definecolor{noteBg}{RGB}{255,247,224}
\definecolor{noteBorder}{RGB}{209,163,45}

\lstset{
  backgroundcolor=\color{codeBg},
  basicstyle=\ttfamily\small,
  breaklines=true,
  frame=single,
  rulecolor=\color{gray!40},
  columns=fullflexible,
  keepspaces=true,
  showstringspaces=false,
  xleftmargin=6pt,
  framexleftmargin=6pt
}

% simple "NOTE" callout
\newcommand{\note}[1]{%
  \par\vspace{4pt}%
  \noindent\fcolorbox{noteBorder}{noteBg}{%
    \parbox{\dimexpr\linewidth-2\fboxsep-2\fboxrule}{\small\textbf{Note.}~#1}}%
  \par\vspace{4pt}%
}

% figure helper: \fig{file}{caption}{label}{width-fraction}
\newcommand{\fig}[4]{%
  \begin{figure}[H]\centering
    \includegraphics[width=#4\linewidth]{#1}
    \caption{#2}\label{#3}
  \end{figure}}

\captionsetup{font=small,labelfont=bf}

\title{\textbf{Creating a Niagara N4 Module}\\[2pt]
\large Workbench \emph{New Module} Wizard + IntelliJ IDEA project setup}
\author{niagara-research — operator workflow manual\\\small Worked example: \texttt{ColdRoomPan} (cold-room control module)}
\date{}

\begin{document}
\maketitle
\tableofcontents
\bigskip
\hrule
\bigskip

% =====================================================================
\section{Scope and prerequisites}
% =====================================================================
This manual documents, step by step, how an operator creates the empty
\textbf{skeleton} of a Niagara~N4 module with the Workbench \emph{New Module}
wizard, and then opens and configures that project in \textbf{IntelliJ IDEA}
so it is ready to program. It is the ``manual'' path — the wizard generates
a standard \texttt{gradle-niagara} project that you afterwards edit in the IDE.

The worked example throughout is the module used for the cold-room logic:
\texttt{ColdRoomPan} (symbol \texttt{CRP}, vendor \texttt{Angeles}).

\paragraph{Prerequisites.}
\begin{itemize}[nosep]
  \item Niagara Workbench (here: \emph{Optimizer Supervisor} N4.14.0.162).
  \item IntelliJ IDEA (here: 2026.1.1).
  \item A full \textbf{JDK~8} (here: Azul \emph{Zulu}~8) — Niagara builds require
        a real JDK, not just a JRE.
  \item Optionally Node.js, only if the module will build JavaScript (UX) parts.
\end{itemize}

\note{This produces exactly the kind of \texttt{gradle-niagara} project that the
repository's \texttt{docs/module-dev-workflow.md} runbook describes; once it is
open and configured, you continue with the \emph{edit\,$\rightarrow$\,build\,$\rightarrow$\,sign\,$\rightarrow$\,deploy\,$\rightarrow$\,test} loop.}

% =====================================================================
\section{Part A — Create the module in Workbench}
% =====================================================================

\subsection{Open the New Module wizard}
In Workbench, open the \textbf{Tools} menu and choose \textbf{New Module}
(Figure~\ref{fig:tools}).

\fig{img04.png}{Workbench \texttt{Tools $\rightarrow$ New Module}.}{fig:tools}{0.72}

\subsection{Step 1 of 3 — module metadata and runtime profiles}
The wizard opens on \emph{Step 1 of 3}. The first field is the
\textbf{Create Module under Directory} path — the folder where the project
will be generated (Figure~\ref{fig:step1empty}). In this example:
\path{C:/modulos_niagara_n4/Cliente/Leon-Guanjuato/Paccadia}.

\fig{img05.png}{Step 1 of 3, empty. The directory path is selected first.}{fig:step1empty}{0.62}

Fill in the remaining fields (Figure~\ref{fig:step1all}):

\begin{itemize}[nosep]
  \item \textbf{Module Name} — \texttt{ColdRoomPan}. The module's technical name.
  \item \textbf{Preferred Symbol} — \texttt{CRP}. Short prefix/symbol for the module.
  \item \textbf{Version} — \texttt{1.0}.
  \item \textbf{Description} — \texttt{Logic Cold Room Panccadia}.
  \item \textbf{Vendor} — \texttt{Angeles}.
  \item \textbf{Runtime Profiles} — which JAR profiles the module ships (see below).
  \item \textbf{Create Palette} — when checked, generates a palette file so the
        module's components can be dropped from the Workbench palette.
\end{itemize}

\fig{img06.png}{Step 1 filled in, with RUNTIME + UX + WB profiles and
\emph{Create Palette} checked (first attempt).}{fig:step1all}{0.62}

\paragraph{Runtime Profiles — what each one means.}
\begin{description}[nosep,leftmargin=1.6cm,style=nextline]
  \item[RUNTIME (rt)] Core runtime Java classes only, no user interface.
  \item[UX] Lightweight HTML5 + JavaScript + CSS user interface.
  \item[WB] Workbench (or Workbench applet) user-interface classes.
  \item[SE] Java classes that use the full Java~8 Standard Edition platform API.
\end{description}

\note{The recommended choice for this project is \textbf{pure RUNTIME (rt)}.
The \emph{UX} and \emph{WB} profiles are only needed if you want a browser UI or a
Workbench view inside the module — not required for control logic. Starting
RT-only keeps the module small and simple; UX/WB can be added later if a
Workbench interface is actually wanted.}

\subsection{Step 2 of 3 — dependencies}
Clicking \textbf{Next} advances to \emph{Step 2 of 3}, \textbf{Select
Dependencies}. By default the module depends on \texttt{nre} and \texttt{baja}
(both Tridium~4.14), shown in Figure~\ref{fig:step2}. Use \textbf{Add} to pull in
more modules (e.g. \texttt{control}, \texttt{kitControl}, \texttt{alarm},
\texttt{history}, \texttt{schedule}) from the picker in Figure~\ref{fig:step2add}.

\fig{img07.png}{Step 2 of 3 — default dependencies \texttt{nre} and \texttt{baja}.}{fig:step2}{0.62}
\fig{img08.png}{The \emph{Select Dependencies} picker (Add).}{fig:step2add}{0.7}

\note{Dependencies can also be added \emph{later}, directly in the module's build
configuration, so it is fine to leave the defaults here and add what the logic
needs during development.}

\subsection{Step 3 of 3 — packages}
\textbf{Next} advances to \emph{Step 3 of 3}, \textbf{Add Packages}. The wizard
proposes one Java package per selected runtime profile
(Figure~\ref{fig:step3all}): \texttt{com.angeles.ColdRoomPan} (rt),
\texttt{com.angeles.ColdRoomPan.ui} (wb) and \texttt{com.angeles.ColdRoomPan.ux}
(ux). The package prefix follows the vendor (\texttt{com.angeles}).

\fig{img09.png}{Step 3 of 3 — one package per profile (rt/wb/ux) in the
first attempt.}{fig:step3all}{0.62}

\subsection{Gotcha: you cannot go Back to Step~1}
Once you press \textbf{Next} past Step~1, the \textbf{Back} button no longer
returns to Step~1. To change the runtime profiles (here: to drop UX/WB and keep
only RT) you must \textbf{cancel and start the wizard again}.

The second run keeps only \textbf{RUNTIME} plus \emph{Create Palette}
(Figure~\ref{fig:step1rt}), which yields a single package
\texttt{com.angeles.ColdRoomPan} with profile \texttt{rt}
(Figure~\ref{fig:step3rt}).

\fig{img10.png}{Second run — Step 1 with only RUNTIME + Create Palette.}{fig:step1rt}{0.62}
\fig{img11.png}{Second run — Step 3 with a single \texttt{rt} package.}{fig:step3rt}{0.62}

\subsection{Finish — the generated project}
Pressing \textbf{Finish} creates the project. Browsing the target folder in
Workbench shows the generated skeleton (Figure~\ref{fig:files}):

\begin{itemize}[nosep]
  \item \texttt{ColdRoomPan/} — the module source directory.
  \item \texttt{gradle/} — the Gradle wrapper support files.
  \item \texttt{build.gradle.kts}, \texttt{settings.gradle.kts} — Gradle build scripts.
  \item \texttt{gradle.properties} — build configuration (paths, JDK) — edited later.
  \item \texttt{gradlew}, \texttt{gradlew.bat} — the Gradle wrapper launchers.
\end{itemize}

\fig{img12.png}{The generated project files in the target directory
(\texttt{Paccadia}).}{fig:files}{0.95}

At this point the skeleton exists; what remains is to open it in the IDE and
point the build at the correct Niagara install and JDK.

% =====================================================================
\section{Part B — Open the project in IntelliJ IDEA}
% =====================================================================

\subsection{Launch IntelliJ IDEA}
Open IntelliJ IDEA (Figure~\ref{fig:ijsearch}).

\fig{img13.png}{Launching IntelliJ IDEA 2026.1.1.}{fig:ijsearch}{0.7}

\subsection{Open the module folder}
On the Welcome screen (Figures~\ref{fig:ijwelcome} and~\ref{fig:ijopenbtn}) click
\textbf{Open}, then browse to and select the module folder \texttt{Paccadia}
(Figure~\ref{fig:ijopen}) and press \textbf{Select Folder}.

\fig{img14.png}{IntelliJ Welcome screen — recent projects.}{fig:ijwelcome}{0.7}
\fig{img15.png}{Welcome screen — \emph{New Project / Open / Clone Repository}.}{fig:ijopenbtn}{0.7}
\fig{img16.png}{\emph{Open File or Project} — select the \texttt{Paccadia} folder.}{fig:ijopen}{0.8}
\fig{img17.png}{Clicking \textbf{Open} to browse for the project.}{fig:ijopen2}{0.7}

\subsection{Trust the project}
IntelliJ asks whether to trust the folder (Figure~\ref{fig:trust}). Choose
\textbf{Trust Project}. (The dialog optionally adds the IDE and project folders
to the Microsoft Defender exclusion list to improve performance.)

\fig{img19.png}{\emph{Trust and Open Project 'Paccadia'?} $\rightarrow$ Trust Project.}{fig:trust}{0.7}

After trusting, IntelliJ begins importing the Gradle project
(Figure~\ref{fig:import}); wait for the import to finish.

\fig{img18.png}{IntelliJ importing the \texttt{Paccadia} Gradle project.}{fig:import}{0.95}

% =====================================================================
\section{Part C — Configure the project}
% =====================================================================

\subsection{Set the SDK and language level}
Open \textbf{File $\rightarrow$ Project Structure} (Figure~\ref{fig:psmenu}).

\fig{img20.png}{\texttt{File $\rightarrow$ Project Structure}.}{fig:psmenu}{0.75}

Under \textbf{Project Settings $\rightarrow$ Project}
(Figure~\ref{fig:psdialog}) set the two values Niagara~N4 requires:
\begin{itemize}[nosep]
  \item \textbf{SDK}: \texttt{1.8} (here \texttt{1.8.0\_472}).
  \item \textbf{Language level}: \texttt{8 - Lambdas, type annotations, etc.}
\end{itemize}
Apply and wait a moment for the changes to take effect.

\fig{img21.png}{Project Structure — SDK 1.8 and language level 8.}{fig:psdialog}{0.8}

\subsection{Edit \texttt{gradle.properties}}
Open \texttt{gradle.properties} (Figure~\ref{fig:gpbefore}). This file tells the
build which Niagara install and which JDK to use.

\fig{img22.png}{\texttt{gradle.properties} before editing (JDK lines commented out).}{fig:gpbefore}{0.95}

The relevant keys, as generated:
\begin{lstlisting}
# The path to the installation of Niagara you are building against
niagara_home=C:\\Honeywell\\OptimizerSupervisor-N4.14.0.162

# The path to niagara_user_home for the version you are building against
niagara_user_home=C:\\Users\\equipo\\Niagara4.14\\OptimizerSupervisor

# Uncomment and set to the path of your node install if building JavaScript
#nodeHome=

#org.gradle.java.installations.auto-detect=false
#org.gradle.java.installations.auto-download=false
#org.gradle.java.installations.paths=
\end{lstlisting}

Apply the edits (uncomment and set the JDK toolchain, and Node if used):
\begin{lstlisting}
niagara_home=C:\\Honeywell\\OptimizerSupervisor-N4.14.0.162
niagara_user_home=C:\\Users\\equipo\\Niagara4.14\\OptimizerSupervisor

nodeHome=C:\\Program Files\\nodejs

org.gradle.java.installations.auto-detect=false
org.gradle.java.installations.auto-download=false
org.gradle.java.installations.paths=C:\\Program Files\\Zulu\\zulu-8
\end{lstlisting}

\note{Pinning \texttt{org.gradle.java.installations.paths} to the Zulu~8 JDK and
turning \emph{auto-detect} / \emph{auto-download} off guarantees the build uses
the correct Java~8 toolchain instead of whatever Gradle might pick.}

\subsection{The two version-dependent edits}
When you target a \emph{different} Niagara version, two things change:
\begin{enumerate}[nosep]
  \item \texttt{gradle.properties} — update \texttt{niagara\_home} (and
        \texttt{niagara\_user\_home}) to the path of that Niagara install.
  \item \texttt{settings.gradle.kts} — update the Gradle version to the one that
        Niagara version builds with.
\end{enumerate}

\subsection{Result — environment ready}
Once the Gradle project finishes importing with the correct SDK and paths, the
Gradle tool window shows the project and its \texttt{ColdRoomPan-rt} sub-project
with its tasks (Figure~\ref{fig:gradleloaded}). The environment is now ready to
start programming the module.

\fig{img24.png}{Gradle tool window — \texttt{ColdRoomPan} and
\texttt{ColdRoomPan-rt} loaded.}{fig:gradleloaded}{0.85}

% =====================================================================
\section{Part D — Programming: Gradle tasks and key files}
% =====================================================================
Once the project is open and configured you program the module and drive the
build from the \textbf{Gradle} tool window. Remember that the \textbf{Gradle
version depends on the Niagara version} (set in \texttt{settings.gradle.kts}).

\subsection{Open the Gradle tool window (elephant icon)}
Open the Gradle panel from the elephant icon on the right edge
(Figures~\ref{fig:gradleopen} and~\ref{fig:gradleroot}). It lists the root
project \texttt{ColdRoomPan}, its \textbf{Tasks}, and the \texttt{ColdRoomPan-rt}
sub-project.

\fig{img23.png}{Opening the Gradle tool window.}{fig:gradleopen}{0.72}
\fig{img24.png}{Gradle root project \texttt{ColdRoomPan} with its \emph{Tasks}
and the \texttt{ColdRoomPan-rt} part.}{fig:gradleroot}{0.85}

\subsection{The tasks that matter: build, clean, slotomatic}
Expand \textbf{Tasks} (Figure~\ref{fig:tasks}). Two groups matter:
\begin{itemize}[nosep]
  \item \textbf{build} — \texttt{assemble}, \texttt{build},
        \texttt{buildDependents}, \texttt{buildNeeded}, \texttt{classes},
        \texttt{clean}, \texttt{jar}, \texttt{moduleTestClasses},
        \texttt{moduleTestJar}, \texttt{testClasses}. The two you use most are
        \texttt{build} and \texttt{clean}.
  \item \textbf{niagara} — \texttt{migrateSlotomatic} and \texttt{slotomatic}.
        \texttt{slotomatic} is the code-generation task.
\end{itemize}

\fig{img25.png}{Tasks expanded: the \emph{build} gear tasks and the
\emph{niagara} group with \texttt{slotomatic}.}{fig:tasks}{0.62}

\note{\textbf{Slotomatic rule} (documented in \texttt{docs/module-dev-workflow.md}
§1.3/§3 and \texttt{docs/module-best-practices.md} §5.2): run \texttt{:slotomatic}
\emph{only} when a \texttt{@Niagara*} annotation changed. It writes the AUTO slot
region of the class; \textbf{never} hand-edit inside the AUTO markers. A stale
AUTO region causes a compile/runtime failure. Typical loop:
\emph{clean $\rightarrow$ slotomatic (only on annotation change) $\rightarrow$ build}.}

\subsection{Key module files to keep in mind}
The \texttt{ColdRoomPan-rt} part contains the Niagara descriptor files
(Figure~\ref{fig:keyfiles}):
\begin{description}[nosep,leftmargin=0.5cm,style=nextline]
  \item[\texttt{module.palette}] BOG of components for drag-and-drop from the
    Workbench palette. Ship it for any module that exports reusable components
    (ungated, lazy, fault-tolerant). Documented: \texttt{module-best-practices.md} [B634].
  \item[\texttt{module-permissions.xml}] A \emph{Java Security Manager
    permission-request manifest}. It declares which restricted host capabilities
    the module's own \emph{code} needs granted — network sockets, file access,
    native libraries, authentication, system time, backups, keystore — so the
    module can run protected operations without an \texttt{AccessControlException}.
    Root \texttt{<permissions>}, with \texttt{<niagara-permission-groups
    type="station|workbench|all">} blocks holding \texttt{<req-permission>} /
    \texttt{<opt-permission>} entries (\texttt{<name>}, \texttt{<purposeKey>},
    optional \texttt{<parameters>} such as hosts/ports). It does \textbf{not}
    control how the module looks, and it is \textbf{not} per-user RBAC. The
    devkit emits it as a fully-commented stub (no extra permissions requested);
    the \texttt{chihuahua} module ships exactly that stub. Doc:
    devguide \texttt{security/requestingPermissions.txt} §1, §2.1.1, §2.2;
    build ref \texttt{build.txt:160}.
  \item[\texttt{module.lexicon}] Localization strings for the module.
  \item[\texttt{niagara-module.xml} / \texttt{module-include.xml}] The exported
    \texttt{<type>} list that Slotomatic reads.
  \item[\texttt{ColdRoomPan-rt.gradle.kts}] The part build script:
    \texttt{moduleManifest} (\texttt{moduleName}, \texttt{runtimeProfile=rt}),
    \texttt{dependencies} (\texttt{nre(":nre")}, \texttt{api(":baja")},
    \texttt{moduleTestImplementation(":test-wb")}), and the \texttt{bajadoc} task
    (\texttt{includePackage("com.angeles.ColdRoomPan")}).
\end{description}

\fig{img26.png}{The \texttt{ColdRoomPan-rt} part: \texttt{module.palette},
\texttt{module-permissions.xml}, \texttt{module.lexicon},
\texttt{niagara-module.xml}, and \texttt{ColdRoomPan-rt.gradle.kts}.}{fig:keyfiles}{0.95}

\subsection{Reference material}
\begin{itemize}[nosep]
  \item Local tooling: \path{\\wsl.localhost\Ubuntu\home\cristian\modulos_niagara_n4\niagara-tools}
  \item Real worked example to study a full module's structure:
        \path{\\wsl.localhost\Ubuntu\home\cristian\modulos_niagara_n4\Cliente\Honeywell\MX60\chihuahua}
  \item Repository runbooks that document this part:
        \texttt{docs/module-dev-workflow.md} and \texttt{docs/module-best-practices.md}.
\end{itemize}

% =====================================================================
\section{Part E — Building from WSL2 (command line)}
% =====================================================================
The project is created on Windows (the wizard) but can be developed and built
from a WSL2 (Ubuntu) shell — the pattern used for the \texttt{chihuahua} module.
This section documents that build loop. Source of truth: the \texttt{chihuahua}
project's \texttt{BUILD\_WORKFLOW.md} / \texttt{CLAUDE.md}.

\subsection{Prerequisites in WSL2}
\begin{itemize}[nosep]
  \item A \textbf{JDK 8} installed in WSL: \texttt{sudo apt-get install -y openjdk-8-jdk}
        (lands at \path{/usr/lib/jvm/java-8-openjdk-amd64}).
  \item The Niagara install reachable from WSL at
        \path{/mnt/c/Honeywell/OptimizerSupervisor-N4.14.0.162}.
\end{itemize}

\subsection{Mandatory \texttt{-P} overrides}
\texttt{gradle.properties} points at Windows \texttt{C:\textbackslash} paths that
do not exist inside WSL, so every WSL build must override them on the command line:
\begin{lstlisting}
-Pniagara_home=/mnt/c/Honeywell/OptimizerSupervisor-N4.14.0.162
-Porg.gradle.java.installations.paths=/usr/lib/jvm/java-8-openjdk-amd64
\end{lstlisting}

\subsection{The build loop (from the project root)}
\begin{lstlisting}
# Clean + (regenerate AUTO region) + build the signed jar
./gradlew :ColdRoomPan-rt:clean :ColdRoomPan-rt:slotomatic :ColdRoomPan-rt:jar \
  -Pniagara_home=/mnt/c/Honeywell/OptimizerSupervisor-N4.14.0.162 \
  -Porg.gradle.java.installations.paths=/usr/lib/jvm/java-8-openjdk-amd64
\end{lstlisting}
\begin{itemize}[nosep]
  \item \texttt{slotomatic} runs in WSL and regenerates the AUTO GENERATED region
        from the \texttt{@Niagara*} annotations. Run it \emph{only} when an
        annotation changed; otherwise \texttt{clean + jar} is enough. ``Cannot
        transform \dots\ no metadata'' warnings for \texttt{srcTest/} are benign.
  \item The signed jar is produced by the \texttt{niagara-signing} plugin and lands
        in \path{$niagara_home/modules/ColdRoomPan-rt.jar}; copy it to the target
        station's \texttt{modules} directory to deploy.
  \item \textbf{Version bump}: when you add or change a \emph{frozen} slot
        (\texttt{@NiagaraProperty} / \texttt{@NiagaraAction}), bump
        \texttt{defaultModuleVersion} in the root \texttt{build.gradle.kts}
        (chihuahua's \texttt{deploy.sh --bump}), or the new slot will not appear.
\end{itemize}

\subsection{Quick compile-only check}
Before a full build you can validate that the sources compile (useful while
integrating code):
\begin{lstlisting}
./gradlew :ColdRoomPan-rt:compileModuleTestJava \
  -Pniagara_home=/mnt/c/Honeywell/OptimizerSupervisor-N4.14.0.162 \
  -Porg.gradle.java.installations.paths=/usr/lib/jvm/java-8-openjdk-amd64
\end{lstlisting}

\note{Signing: the \texttt{niagara-signing} plugin signs the jar using the
certificate store in \texttt{niagara\_user\_home}. If signing fails under WSL, also
override it: \path{-Pniagara_user_home=/mnt/c/Users/equipo/Niagara4.14/OptimizerSupervisor}.
Deploying a custom module to a locked-down OEM (Honeywell) install requires a
certificate the station trusts (here: the \texttt{ANGELES} identity).}

\subsection{Tests}
Do \textbf{not} gate the build on \texttt{gradle test} / \texttt{niagaraTest}: the
\texttt{com.tridium.niagara-module} 7.6.17 plugin has a known bug
(\texttt{moduleTestAnnotationProcessor} produces no metadata, so \texttt{@NiagaraType}
tests are silently skipped, \texttt{Total tests run: 0}). Test files in
\texttt{srcTest/} are documentation-only skeletons.

\subsection{Operator-only activation steps}
The build/deploy runs from WSL, but these are done by the operator on the station:
\begin{itemize}[nosep]
  \item \textbf{Restart the station} — for Java changes and frozen slots.
  \item \textbf{Close and reopen Workbench} — if a frozen slot/action still does not
        appear after the station restart (Workbench's client type cache does not
        refresh on a station restart).
  \item \textbf{Smoke test} — verify on the live station.
\end{itemize}

\note{Gotcha: \texttt{compileModuleTestJava} pollutes \texttt{module-include.xml}
with test types — \texttt{git checkout} that file before committing.}

% =====================================================================
\section{Summary}
% =====================================================================
\begin{enumerate}[nosep]
  \item Workbench: \texttt{Tools $\rightarrow$ New Module}.
  \item Step~1: path + Module Name, Preferred Symbol, Version, Description,
        Vendor; choose \textbf{RUNTIME} only; check \textbf{Create Palette}.
        (Remember: no going Back to Step~1.)
  \item Step~2: dependencies (defaults \texttt{nre}, \texttt{baja}; add more here
        or later).
  \item Step~3: one \texttt{rt} package \texttt{com.<vendor>.<Module>}; Finish.
  \item IntelliJ: \textbf{Open} the generated folder $\rightarrow$
        \textbf{Trust Project} $\rightarrow$ wait for the Gradle import.
  \item \textbf{Project Structure}: SDK \texttt{1.8}, language level \texttt{8}.
  \item \texttt{gradle.properties}: set \texttt{niagara\_home},
        \texttt{niagara\_user\_home}, \texttt{nodeHome} (if needed), and pin the
        Zulu~8 JDK toolchain.
  \item Per Niagara version: adjust \texttt{niagara\_home} in
        \texttt{gradle.properties} and the Gradle version in
        \texttt{settings.gradle.kts}.
\end{enumerate}

\paragraph{Example values used in this manual.}
\begin{center}
\begin{tabular}{@{}ll@{}}
\toprule
Field & Value \\
\midrule
Module Name & \texttt{ColdRoomPan} \\
Preferred Symbol & \texttt{CRP} \\
Version & \texttt{1.0} \\
Description & \texttt{Logic Cold Room Panccadia} \\
Vendor & \texttt{Angeles} \\
Package (rt) & \texttt{com.angeles.ColdRoomPan} \\
Directory & \path{C:/modulos_niagara_n4/Cliente/Leon-Guanjuato/Paccadia} \\
Niagara home & \path{C:/Honeywell/OptimizerSupervisor-N4.14.0.162} \\
JDK & Zulu 8 (\path{C:/Program Files/Zulu/zulu-8}) \\
\bottomrule
\end{tabular}
\end{center}

\end{document}
